\section{Model}

\subsection{Model of Problem 1}

\paragraph{Force analysis.} The body is acted on by three forces: the lateral disturbance
$F_d(t)$, the spring--damper constraint force, and the actuator resultant force. Taking
the lateral direction as positive,
\begin{equation}
  F_h(t) = -k\,y(t) - c\,\dot y(t),
  \label{eq:fh}
\end{equation}
with $k$ and $c$ listed in Table~\ref{tab:params}.

\paragraph{Actuator force.} A single eccentric block of radius $r$ rotating at
$\omega(t)$ exerts $m r \omega^2$ along its radius. The actuator carries two blocks
rotating in opposite directions at equal speed, so the lateral components add and the
in-plane components cancel; for two actuators the resultant is
\begin{equation}
  F_a(t) = 2mr\left[\omega_i^2\cos\theta_i + \alpha_i\sin\theta_i\right] \quad\text{(per actuator pair)}.
  \label{eq:fa}
\end{equation}

\paragraph{Equation of motion.} Substituting \eqref{eq:fh} and \eqref{eq:fa} into Newton's
second law for the lateral translation gives
\begin{equation}
  M\ddot y(t) + c\dot y(t) + k y(t) = F_d(t) + F_a(t),
  \label{eq:eom}
\end{equation}
which is integrated with a fourth-order Runge--Kutta scheme at $\Delta t$, with the
angular velocity and acceleration obtained from $\theta_i(t)$ by finite differences.

\paragraph{Vibration index.} With $T = 10$\,s,
\begin{equation}
  I_h = \frac{1}{T}\int_0^{T} a_y^2(t)\,\mathrm{d}t,
  \label{eq:index}
\end{equation}
evaluated by trapezoidal integration.

% 参数表：一次写清，后面正文不再重复数字。
\threelinetable{Model parameters.}{tab:params}{@{}p{0.34\textwidth}lrr@{}}
  {Quantity & Symbol & Value & Unit}
  {Body mass & $M$ & 2000 & kg \\
   Eccentric mass & $m$ & 100 & kg \\
   Eccentric radius & $r$ & 0.2 & m \\
   Lateral stiffness & $k$ & $7.2252\times10^{6}$ & N/m \\
   Lateral damping & $c$ & 3600 & N$\cdot$s/m \\
   Angular velocity limit & $\omega_{\max}$ & 100 & rad/s \\
   Angular acceleration limit & $\alpha_{\max}$ & 5000 & rad/s$^2$}

\subsection{Result of Problem 1}

% 结果先核约束、再报指标——这是原稿做得最好的一个动作，务必保留。
% 只报"降了 80 %"却不证明作动器没超限，评委会怀疑方案不可执行。
With the actuator inactive the index is $I_h^{\text{off}} = [\text{value}]$ [unit]; with the
actuator driven by the prescribed schedule it is $I_h^{\text{on}} = [\text{value}]$ [unit].
Table~\ref{tab:limits} confirms the schedule stays inside the actuator's envelope, so the
comparison is meaningful.

\begin{table}[htbp]
  \centering
  \caption{Constraint check and index summary for the three problems.}
  \label{tab:limits}
  \small
  \begin{tabular}{@{}lrrrr@{}}
    \toprule
    Case & $I_h$ off & $I_h$ on & Peak $\omega$ (rad/s) & Peak $\alpha$ (rad/s$^2$) \\
    \midrule
    Problem 1 & [ ] & [ ] & [ ] & [ ] \\
    Problem 2 & [ ] & [ ] & [ ] & [ ] \\
    Problem 3 & [ ] & [ ] & [ ] & [ ] \\
    \midrule
    Limit     & --  & --  & 100 & 5000 \\
    \bottomrule
  \end{tabular}
\end{table}

\subsection{Model of Problem 2}

% 与问题一的差别只写"差在哪"，不重抄力分析。原稿这一点做对了。
The force analysis is unchanged; the difference is that $\theta_i(t)$ is now the decision
variable. Over a horizon of $N_p$ steps, at each step solve
\begin{equation}
  \min_{\{\theta_i\}} \; \hat I_h
  \quad\text{s.t.}\quad
  M\ddot y + c\dot y + k y = F_d + F_a(\theta,\omega,\alpha),\;\;
  |\omega|\le \omega_{\max},\;\; |\alpha|\le \alpha_{\max},
\end{equation}
apply the first step, and repeat.

\subsection{Result of Problem 2}

Report the index for both states, the suppression rate, the schedule itself (Figure), and
the peak $\omega$ / $\alpha$ against their limits.

\subsection{Model of Problem 3}

% 问题三的重点是"少了什么信息、用什么补回来"，不是再推一遍动力学。
With $F_d$ unmeasurable, only $a_y$ is observed. An extended state observer reconstructs
$y$, $\dot y$ and the lumped disturbance $z$:
\begin{equation}
  \dot{\hat y} = \hat v + \beta_0 (y - \hat y),\quad
  \dot{\hat v} = a_y + \beta_1 (y - \hat y),\quad
  \dot{\hat z} = \beta_2 (y - \hat y),
\end{equation}
and a PID loop on the reconstructed state drives the angular displacement:
\begin{equation}
  \theta(t) = K_p e + K_i\!\int e\,\mathrm{d}t + K_d \dot e .
\end{equation}

\subsection{Result of Problem 3}

Report the index for both states, the suppression rate, the disturbance estimation error
after convergence, and the gain sensitivity (Section 5.3).

% ┈┈┈┈┈┈┈┈┈┈┈┈┈┈┈┈┈┈┈┈┈┈┈┈┈┈┈┈┈┈┈┈┈┈┈
% 【Section writing prompt】
% ════════════════════════════════════════════
% Position: the body of the paper; it should take about 60 % of the page budget.
% Five beats per question, identical across questions: force analysis → equations
%           (numbered) → numerical method → results → index.
% Results:  check the actuator limits BEFORE quoting the index. A suppression rate with
%           no feasibility check invites the reader to doubt the scheme.
% Tables:   at least three — parameters, index summary, constraint check. A paper whose
%           results exist only as curves cannot be audited.
% Do not:   re-derive what Problem 1 already derived. State the difference and move on.
% Units:    one unit per quantity for the whole paper; the index unit must match the
%           dimension of the defining equation, not a plausible-looking one.
% ════════════════════════════════════════════
